\chapter{Conclusion and Future Work}
The internship investigated free-acid monitoring for a TSP process using one month of one-minute production data. The project moved from exploration and feature engineering to controlled benchmarks, target-anchored delta forecasting, target-free virtual sensing, stability research, dashboards and a Docker/OPC UA laboratory prototype.

The main forecasting result is a hybrid-feature Ridge delta model that beat persistence at 1, 5, 10 and 15 minutes on the held-out January period. A separate target-free Ridge-LightGBM virtual sensor outperformed a constant baseline, while five-cluster Ridge was retained by the deployment runtime as a stability-oriented candidate after chronological-fold analysis. The report does not combine these results: they answer different questions and have different evidence levels.

Future work should first obtain a later labeled production period and evaluate frozen candidates without retuning. It should then perform rolling-origin multi-month validation, verify units and proxy assumptions with process engineers, define laboratory freshness and warning costs, calibrate drift and uncertainty on independent residuals, and validate clusters against operational records. On the deployment side, the next milestones are secured Kepware integration, historian backfill, retention and backup policy, longer read-only shadow operation and operator feedback. Any move beyond advisory use would require separate governance and control-system approval.
